% !TeX root = ../main.tex
\begin{abstract}
MXFP8和MXFP4等微缩格式在大型语言模型中得到了更广泛的应用。这一趋势推动了对基于点积的张量核中增加低精度算力的需求。然而，不同多精度算力配比的硬件成本和工作负载效益缺乏系统性的评估。为了填补这一空白，我们提出了DP-Puzzle，这是一个基于可组合点积单元的设计空间探索框架。这些单元在不同精度之间共享计算数据路径，并包含专用的低精度扩展以支持不同的峰值算力配比。该框架将电路级面积和功耗特性与当前工作负载和潜在精度偏移的模型相结合。在计算核心的总面积预算固定的情况下，DP-Puzzle评估了候选设计的计算延迟和能耗。这一分析为硅前芯片设计期间基于工作负载的算力配比选择提供了量化依据。

\end{abstract}
\begin{IEEEkeywords}
设计空间探索，多精度算力配比，组合式点积单元，张量计算，硬件成本建模
\end{IEEEkeywords}
